\section{Recursive Spectral Zero Hierarchy}
\label{sec:unified-reality-recursive-spectral-zero-hierarchy}

The quartic spectral carrier, finite-field elliptic zeta functions, and
global zeta or $L$-functions share a recursive spectral mechanism.  They
do not share one literal zero set.  The commonality is structural:
finite observations are folded into a generating object, the generating
object exposes spectral factors, normalization identifies the balanced
surface, and zero constraints are read only through the ledgers available
at that level.

\begin{definition}[Zeta-like spectral hierarchy]
\label{def:unified-reality-zeta-like-spectral-hierarchy}
The \emph{zeta-like spectral hierarchy} separates three levels:
$$
\mathsf{RecurrenceSpectralZeros}
\longrightarrow
\mathsf{CurveZetaZeros}
\longrightarrow
\mathsf{GlobalZetaLFunctionZeros}.
$$
The first level contains zeros such as the endpoint zeros of the fixed
quartic recurrence carrier.  The second level contains zeta zeros of
geometric objects over finite fields, such as numerator roots of
elliptic Hasse--Weil zeta functions.  The third level contains zeros of
global arithmetic zeta and $L$-functions, including the Riemann zeta
zero ledger.  The arrows are carrier-role refinements, not inclusions of
one point set into another.
\end{definition}

\begin{definition}[Recursive spectral fold]
\label{def:unified-reality-recursive-spectral-fold}
A \emph{recursive spectral fold} is a packet
$$
\mathsf{ObsTime}
\to
\mathsf{GenFunction}
\to
\mathsf{SpectralFactors}
\to
\mathsf{Normalization}
\to
\mathsf{ZeroConstraint},
$$
where finite observations indexed by an observation-time parameter are
compressed into a generating function or Euler product; spectral factors
read the compressed carrier; normalization identifies the balanced
surface; and zero constraints are derived only after the required
dominance, symmetry, exactness, or continuation ledgers have been
displayed.
\end{definition}

\begin{theorem}[Recursive spectral folds are not self-reference]
\label{thm:unified-reality-recursive-spectral-fold-not-self-reference}
A recursive spectral fold is a hierarchical compression of finite
observation rows, not a circular definition of its own zero set.
\end{theorem}

\begin{proof}
The input rows are finite observations such as recurrence coefficients,
point counts over finite extensions, or prime-local Euler factors.  The
generating object is constructed from those rows.  Spectral factors are
then read from the generated object, and zero constraints are read from
the spectral factors plus the relevant ledgers.  Each stage consumes the
previous stage as data; no stage assumes the zero constraint that it is
supposed to produce.
\end{proof}

\begin{definition}[Quartic recurrence zero level]
\label{def:unified-reality-quartic-recurrence-zero-level}
The quartic recurrence zero level is the recursive spectral fold
$$
(Z_m(y))_{m\ge0}
\to
Z(t,y)=\frac{N(t,y)}{D(t,y)}
\to
\Pi(\lambda,y)
\to
E:\ Y^2=4\lambda^3-4\lambda+1
\to
\mathsf{EndpointZeros}.
$$
Its zeros are zeros of the finite recurrence polynomials $Z_m(y)$ under
the endpoint scaling and fixed-window ledgers of
\autoref{def:unified-reality-quartic-endpoint-zero-reality}.
\end{definition}

\begin{theorem}[Quartic endpoint zeros are recurrence-spectral zeros]
\label{thm:unified-reality-quartic-endpoint-zeros-recurrence-spectral}
The endpoint zeros in the fixed quartic carrier are recurrence-spectral
zeros.  They are zeta-like by recursive spectral form, but they are not
zeros of the Riemann zeta function or of an elliptic Hasse--Weil zeta
function.
\end{theorem}

\begin{proof}
By
\autoref{thm:unified-reality-quartic-spectral-common-carrier-principle},
the fixed quartic carrier transports recurrence algebra through the
spectral polynomial, elliptic normalization, real projection, dominance,
and endpoint-zero ledgers.  The endpoint zeros are obtained by local
branch interference and Rouch\'e isolation in
\autoref{thm:unified-reality-quartic-endpoint-zero-branch-interference}.
No Euler product, finite-field point-count zeta function, or completed
Riemann zeta zero ledger is part of this carrier.  Thus the analogy is
structural rather than an identification of zero sets.
\end{proof}

\begin{definition}[Finite-field elliptic zeta level]
\label{def:unified-reality-finite-field-elliptic-zeta-level}
For an elliptic curve $E/\mathbb F_q$, the finite-field elliptic zeta
level folds the point-count time series
$$
r\longmapsto \#E(\mathbb F_{q^r})
$$
into
$$
Z(E/\mathbb F_q,T)
:=
\exp\left(
\sum_{r\ge1}
\#E(\mathbb F_{q^r})\frac{T^r}{r}
\right).
$$
For an elliptic curve this has the rational form
$$
Z(E/\mathbb F_q,T)
=
\frac{1-aT+qT^2}{(1-T)(1-qT)}
=
\frac{(1-\alpha T)(1-\beta T)}{(1-T)(1-qT)}.
$$
The spectral roots $\alpha,\beta$ are the Frobenius readouts of the
point-count fold.
\end{definition}

\begin{theorem}[Frobenius roots recursively control point counts]
\label{thm:unified-reality-frobenius-roots-control-point-counts}
In the finite-field elliptic zeta level, the spectral roots determine
the point-count observations by
$$
\#E(\mathbb F_{q^r})
=
q^r+1-\alpha^r-\beta^r.
$$
Thus the point-count time series and the Frobenius spectral roots are
two readouts of the same recursive fold.
\end{theorem}

\begin{proof}
The exponential generating definition records all point-count rows.  The
rational form of the elliptic zeta function packages the same data into
the numerator factors $(1-\alpha T)(1-\beta T)$ together with the two
denominator rows $(1-T)$ and $(1-qT)$.  Expanding the logarithmic
derivative of the displayed rational expression recovers the displayed
point-count formula.  Hence the spectral roots compress and read back
the observation-time sequence.
\end{proof}

\begin{definition}[Finite-field elliptic RH-type normalization]
\label{def:unified-reality-finite-field-elliptic-rh-normalization}
The finite-field elliptic RH-type normalization is the row
$$
|\alpha|=|\beta|=\sqrt q,
$$
equivalently
$$
\left|\frac{\alpha}{\sqrt q}\right|
=
\left|\frac{\beta}{\sqrt q}\right|
=1.
$$
It is the finite geometric form of a normalized spectral-root constraint.
\end{definition}

\begin{theorem}[Finite-field elliptic zeros give a closed local model]
\label{thm:unified-reality-finite-field-elliptic-zeros-closed-model}
The finite-field elliptic zeta level is a closed local geometric model of
RH-type normalization, not a proof of the Riemann zeta zero constraint.
\end{theorem}

\begin{proof}
The normalization row of
\autoref{def:unified-reality-finite-field-elliptic-rh-normalization}
places the Frobenius spectral roots on the unit circle after division by
$\sqrt q$.  This is a theorem at the finite-field elliptic level.  The
Riemann zeta zero ledger instead concerns complex zero rows $s$ of the
completed zeta carrier and the fixed-half condition $\Re(s)=1/2$.
Transport from the finite-field elliptic carrier to the completed
Riemann zeta carrier is not supplied by the finite-field theorem.  The
remaining global bridge is the zero-defect rigidity route isolated in
\autoref{sec:unified-reality-rh-spectral-rigidity-route}.
\end{proof}

\begin{definition}[Global elliptic and zeta zero level]
\label{def:unified-reality-global-elliptic-zeta-zero-level}
For an elliptic curve over $\mathbb Q$, the global elliptic $L$-function
level is assembled from local Euler factors.  At a good prime $p$ the
local factor has the form
$$
1-a_p p^{-s}+p^{1-2s}.
$$
The global readout is an Euler product
$$
L(E,s)=\prod_p L_p(E,p^{-s})^{-1}.
$$
The Riemann zeta readout is the simpler global Euler product
$$
\zeta(s)=\prod_p(1-p^{-s})^{-1}.
$$
Both are global arithmetic spectral carriers, with different local
factor data and different proof obligations.
\end{definition}

\begin{theorem}[Global elliptic L-functions are higher arithmetic siblings]
\label{thm:unified-reality-global-elliptic-lfunctions-higher-siblings}
Global elliptic $L$-functions and the Riemann zeta function share the
Euler-product, Mellin, functional-equation, and zero-ledger pattern, but
their zero ledgers are not the same object.
\end{theorem}

\begin{proof}
\autoref{def:unified-reality-elliptic-modular-lfunction-row} records the
elliptic, modular, and $L$-function bridge shape when the required
bridge rows are displayed.  The Riemann zeta carrier has its own
prime-local factors and completed zero ledger, as used in
\autoref{sec:unified-reality-zeta-ontic-projection}.  The two carriers
share the global arithmetic spectral pattern, but a zero of one carrier
is not automatically a zero of the other.  Any transport between their
zero statements requires its own bridge ledger.
\end{proof}

\begin{definition}[Normalized spectral-root analogy]
\label{def:unified-reality-normalized-spectral-root-analogy}
The normalized spectral-root analogy compares the finite-field elliptic
readout
$$
\frac{\alpha}{\sqrt q}
$$
with the centered prime-channel readout
$$
p^{1/2-s}.
$$
Both are unit-circle constraints after normalization:
$$
\left|\frac{\alpha}{\sqrt q}\right|=1,
\qquad
|p^{1/2-s}|=1.
$$
The first is a Frobenius-root statement in a finite-field curve carrier.
The second is a prime-channel reading of the fixed-half surface in the
zeta carrier.
\end{definition}

\begin{theorem}[Unit-circle analogies require carrier separation]
\label{thm:unified-reality-unit-circle-analogies-carrier-separation}
The analogy between $\alpha/\sqrt q$ and $p^{1/2-s}$ identifies a common
normalization pattern, not a common theorem or a common zero set.
\end{theorem}

\begin{proof}
In the finite-field elliptic carrier, the unit-circle statement is the
Frobenius-root normalization row.  In the Riemann zeta carrier,
$|p^{1/2-s}|=1$ is equivalent to the fixed-half surface by
\autoref{thm:unified-reality-prime-laplace-detect-same-fixed-surface}.
The latter is still only a surface readout until a zero row is shown to
land on that surface.  That additional step is the defect-zero
incompatibility or prime-sieve odd exactness boundary of
\autoref{def:unified-reality-rh-defect-zero-incompatibility-route} and
\autoref{def:unified-reality-prime-sieve-odd-exactness}.
\end{proof}

\begin{definition}[Apophatic far end of recursive spectral towers]
\label{def:unified-reality-recursive-spectral-apophatic-far-end}
The far end of a recursive spectral tower is not an internal point of the
tower.  It is the boundary role named by the continuing availability of
finite observation layers:
$$
r=1,2,3,\ldots
$$
or by unbounded prime windows, recurrence lengths, or extension degrees.
In BEDC notation this is an apophatic far-end role, not a term to be
substituted into the carrier as a value.
\end{definition}

\begin{theorem}[Infinite observation time is not a spectral point]
\label{thm:unified-reality-infinite-observation-time-not-spectral-point}
The infinite direction in a zeta-like recursive spectral fold is a
boundary role for all finite readouts, not a point of the generated
spectral carrier.
\end{theorem}

\begin{proof}
For finite-field elliptic zeta functions, each $r$ supplies the finite
observation $\#E(\mathbb F_{q^r})$.  The zeta function folds the entire
finite observation family into a spectral object.  It does not require a
point $\mathbb F_{q^\infty}$ inside the carrier.  The same discipline is
the apophatic separation of
\autoref{def:unified-reality-cataphatic-process-apophatic-far-end} and
\autoref{prin:apophatic-fixed-point-T-not-kernel}: the far end may be
named as a boundary role, but it is not consumed as a kernel object or
as an additional spectral point.
\end{proof}

\begin{theorem}[Recursive spectral zero hierarchy principle]
\label{thm:unified-reality-recursive-spectral-zero-hierarchy-principle}
Quartic endpoint zeros, finite-field elliptic zeta zeros, global elliptic
$L$-function zeros, and Riemann zeta zeros are projections of a common
recursive spectral pattern only after their carrier boundaries are kept
distinct.
\end{theorem}

\begin{proof}
The quartic carrier uses recurrence data, a rational generating
function, a spectral polynomial, elliptic normalization, dominance, and
endpoint-zero ledgers.  The finite-field elliptic carrier uses point
counts over finite extensions, a Hasse--Weil zeta function, Frobenius
spectral roots, and unit-circle normalization.  The global elliptic and
Riemann zeta carriers use Euler products, completion, functional
equations, and zero ledgers.  These are instances of
\autoref{def:unified-reality-recursive-spectral-fold}.  Their shared
pattern licenses comparison of roles, while their separate ledgers
forbid collapsing one zero reality into another.
\end{proof}
